\section{Construction speed vs.\ a Merkle tree}
\label{sec:benchmark}

Proof size is only half the comparison against Merkle trees; the other is
build time. We measured this directly, using the identical hand-rolled
SHA-256 compression function of the reference implementation
(\texttt{sha256\_core.py}) for every hash performed on every side --
same language, same implementation, no library shortcut on one side and not
the other -- across three constructions:

\begin{enumerate}
\item \textbf{FlatMerkle}, as described in this paper: $2N-2$ items plus one
  finalization, $2N-1$ compression calls total.
\item \textbf{Merkle tree, same convention}: one compression call per node
  (leaf $= \C(\mathrm{IV}, \text{element})$, internal $= \C(\mathrm{IV},
  \text{left} \Vert \text{right})$), no per-node padding -- an
  equally-optimized tree, isolating the tree-vs-chain question with no
  artificial handicap on either side. $N$ leaves $+ (N-1)$ internal nodes
  $= 2N-1$ calls.
\item \textbf{Merkle tree, standard}: every node is its own independently
  padded, finalized SHA-256 call -- i.e.\ what many real Merkle tree
  implementations that hash with a single, ordinary SHA-256 per node
  actually do. A 64-byte node input needs 2 compression calls once padded
  and length-encoded, so this costs $4N-2$ calls total. (Bitcoin's own
  transaction Merkle tree is not an instance of this convention: it
  combines each pair of nodes with \emph{double}-SHA256 -- SHA256 applied
  twice, i.e.\ two independent finalized hashes chained together -- which
  this benchmark does not measure and is more expensive still than either
  convention above.)
\end{enumerate}

\begin{table}[htbp]
\centering
\begin{tabular}{@{}lrrrr@{}}
\toprule
Method & $N$ & Seconds & Compress calls & Calls/element \\
\midrule
FlatMerkle              & 8{,}192  &  1.37 & 16{,}383 & 2.000 \\
Merkle, same convention & 8{,}192  &  1.36 & 16{,}383 & 2.000 \\
Merkle, standard        & 8{,}192  &  2.73 & 32{,}766 & 4.000 \\
\midrule
FlatMerkle              & 65{,}536 & 11.03 & 131{,}071 & 2.000 \\
Merkle, same convention & 65{,}536 & 10.78 & 131{,}071 & 2.000 \\
Merkle, standard        & 65{,}536 & 21.68 & 262{,}142 & 4.000 \\
\bottomrule
\end{tabular}
\caption{Construction speed, real measurements. Calls/element ratios are
the exact rationals $2-1/N$ (FlatMerkle and the same-convention tree) and
$4-2/N$ (the standard tree), which already round to 2.000/4.000 to three
decimals at these scales ($16383/8192 = 1.99988\ldots$) -- not an
empirical coincidence, but not literally exact integers either, only in
the $N \to \infty$ limit.}
\label{tab:benchmark}
\end{table}

Two findings, both borne out at both measured scales:

\begin{enumerate}
\item \textbf{Against an equally-optimized Merkle tree, FlatMerkle is not
  faster to build} -- measured construction time is within 1--2\% either
  way (ratio 0.98--0.99$\times$) at both scales tested, well inside the
  noise expected from a single run each; we do not have repeated-trial
  variance data to call this a rigorously established statistical tie,
  only a close match consistent with the two methods' identical call
  counts. This is expected: FlatMerkle emits $2N-2$ items and a
  same-convention Merkle tree performs $2N-1$ calls -- essentially the same
  work either way. The value of FlatMerkle was never faster construction;
  it is that the resulting commitment is a flat hash of raw bytes
  (Section~\ref{sec:bitcoin-annex}), which a Merkle root is not, and that
  proofs need no separate tree-traversal metadata.
\item \textbf{Against the independently-finalized single-SHA-256
  Merkle-tree convention benchmarked here (\emph{Merkle tree, standard}
  above, matching many real Merkle tree implementations that hash each
  node with one ordinary, self-contained SHA-256 call) -- FlatMerkle is
  almost exactly $2\times$ faster} (ratio 1.96--2.00$\times$), because a
  64-byte node input needs two compression calls once padded and
  length-encoded, where FlatMerkle's continuous stream needs only one call
  per item and pays for padding exactly once, for the whole file. (This
  specific comparison does not include Bitcoin's own transaction Merkle
  tree, which uses double-SHA256 per node rather than this convention --
  see the note above.)
\end{enumerate}

Extrapolating to $N = 2^{20} = 1{,}048{,}576$ elements from the exact
call-count formulas ($2N-1$ for FlatMerkle and the same-convention tree,
$4N-2$ for the standard tree) and the throughput measured at
$N=65{,}536$, pooled across the three methods ($\approx 12{,}057$ calls/s;
the per-method figures, $11{,}880$ / $12{,}164$ / $12{,}093$, differ by
about $\pm2\%$, which we treat as noise) gives $\approx 174$\,s for
FlatMerkle and the same-convention tree and $\approx 348$\,s for the
standard tree. That the first two come out identical is an artifact of
applying one pooled rate to equal call counts, not a claim that equal call
counts guarantee equal wall-clock time. These are computed extrapolations,
not timed runs.

All of these comparisons are single-threaded. FlatMerkle's chaining is
inherently sequential: each checkpoint is the running state after
everything before it. The message-schedule expansion of each of the $N$
data blocks, a substantial part of a compression call, can still be
precomputed in parallel, since data blocks are known in advance; only the
$\mathrm{MID}$ blocks and the chaining itself must run in order. A Merkle
tree's subtrees, by contrast, can be hashed entirely in parallel, so on
multi-core hardware a parallel Merkle build can pull ahead of FlatMerkle
by close to the core count at large $N$. How much of FlatMerkle's
sequential fraction a parallel implementation could hide is not something
we measured here.
